\begin{frame}
	\onslide<1->{
		\frametitle{State Space Models}
		\framesubtitle{Satellite Orbits}
	}
	\begin{itemize}
		\onslide<2->{
			\item 	Simplified model of a circular orbit about a perfectly spherical planet\\  (Note: Earth is an oblate spheroid, not a sphere!)	
		}
		\onslide<3->{
			\item     The state vector consists of position and velocity components in the Planet-Centered Inertial (PCI) frame:

		\[
			\mathbf{x} = 
			\begin{bmatrix}
				x \\ y \\ z \\ u \\ v \\ w
			\end{bmatrix}
		\]


}
	\end{itemize}		
	\onslide<4->{\center --------------------------------------------------------------------}

\end{frame}

\begin{frame}
	\onslide<1->{
	\frametitle{State Space Models}
\framesubtitle{Satellite Orbits}
}
	\onslide<2->{

	The dynamics of the satellite are governed by the following differential equations:
	\begin{align*}
	\dot{x} &= u \\
	\dot{y} &= v \\
	\dot{z} &= w \\
	\end{align*}
}
	\onslide<3->{\center --------------------------------------------------------------------}

\end{frame}

\begin{frame}
	\onslide<1->{
	\frametitle{State Space Models}
\framesubtitle{Satellite Orbits}
}
	\onslide<2->{
	and by:
	\begin{align*}
	\dot{u} &= -\frac{GM}{r^3}x \\
	\dot{v} &= -\frac{GM}{r^3}y \\
	\dot{w} &= -\frac{GM}{r^3}z
	\end{align*}
}
	\onslide<3->{
	Where:
	\begin{itemize}
		\item \(G\) is the gravitational constant.
		\item \(M\) is the mass of the planet.
		\item \(r = \sqrt{x^2 + y^2 + z^2}\) is the distance from the satellite to the center of the planet.
	\end{itemize}
}
	\onslide<4->{\center --------------------------------------------------------------------}

\end{frame}

\begin{frame}
	
	\onslide<1->{
	\frametitle{State Space Models}
\framesubtitle{Satellite Orbits}
	}
	\begin{itemize}
		\onslide<2->{
			\item Note that this is a \textbf{nonlinear} state equation, since the $r$ term is a nonlinear function of $x,y,z$!
		}
		\onslide<3->{
			\item We haven't yet dealt with nonlinear differential equations
		}
		\onslide<4->{
				\item To create a state space model, we must first be \textbf{linearize} the function about an equilibrium (or nominal) point and then use the linearized model in an \textbf{extended} Kalman filter 
		}
		\onslide<5->{
			\item Once the system is linearized, we can identify the states in the usual way
}
		\onslide<6->{
			\item This is beyond the scope of the current video
		}
	\end{itemize}
	\onslide<7->{\center --------------------------------------------------------------------}
	
\end{frame}

